% =========================================================
% 2.4 MICRO BUILT ENVIRONMENT VIA STREET VIEW IMAGERY
% El \section{} de esta sección está en el chapter.tex del capítulo;
% sus referencias van en seccion2_4.bib (misma carpeta).
% =========================================================

En una línea similar de análisis cuantitativo, aunque sin fines de generación de rutas, el siguiente estudio explora la relación entre el crimen y el microentorno construido. Este es un estudio aplicado en 2024 en Toronto mediante imágenes de Google Street View que busca relacionar el crimen con el microentorno construido (MBE, por sus siglas en inglés). Mediante modelos de regresión logística se evaluó la significancia del MBE y de las variables de percepción para clasificar intersecciones de alto y bajo índice delictivo\cite{wu2024analyzing}.

Este análisis permitió asignar importancia a cada uno de los factores considerados, para determinar específicamente cuáles tienen una relación más profunda con el crimen en un segmento determinado, y así identificar qué elementos de la prevención del crimen mediante el diseño urbano deberían tomarse en cuenta en las ciudades\cite{wu2024analyzing}.

Si bien no se desarrolla un aplicativo ni se menciona nada relacionado con rutas, este trabajo tiene una relevancia fundamental para la recolección de información, pues es el único caso de estudio del listado que toma en cuenta tanto datos oficiales de crímenes pasados como elementos de criminología ambiental. Aunado a esto, el análisis se implementa mediante imágenes de Google Street View, lo que empata directamente con la metodología propuesta para el presente proyecto.

El aporte más importante para el proyecto es la serie de variables utilizadas para el estudio de la relación entre el entorno y la criminología, la cual servirá de base para la definición de las propias. Asimismo, la técnica de recolección de imágenes nos brinda un método asequible para el estudio de factores micro y meso en áreas extendidas.
